\begingroup
\setlength{\abovedisplayskip}{7pt}
\setlength{\belowdisplayskip}{7pt}
\setlength{\abovedisplayshortskip}{3pt}
\setlength{\belowdisplayshortskip}{4pt}

% Printed 52 (unfoliated); principal PDF 59; Google PDF 63.
\par\clearpage

\vspace*{.19\textheight}
\begin{center}
{\large IV.}\par\medskip
{\Large BREITENBESTIMMUNG\\DER STERNWARTE SEEBERG.}
\end{center}

Gleichzeitig mit meinen Beobachtungen in Göttingen und Altona wurden
dieselben Sterne auf meine Aufforderung auch von Herrn \textsc{Hansen}, dem Direktor
der Sternwarte Seeberg bei Gotha, mit dem dortigen zweifüßigen \textsc{Ertel}schen
Meridiankreis beobachtet. Der sich daraus ergebende Breitenunterschied zwischen
dieser und der Göttinger Sternwarte ist umso interessanter, als die Seeberger
Sternwarte durch einige unter der Leitung von Generalleutnant \textsc{von Müffling}
gemessene Dreiecke mit dem hannoverschen Dreieckssystem verbunden ist.

Der Kreis wurde während der Beobachtungen einige Male umgelegt. Unabhängig davon
wurde der Kollimationsfehler jeden Tag, meistens zweimal täglich, durch Einstellung
auf den Nadirpunkt bestimmt. Dieses bereits oben erwähnte Verfahren hatte Herr
\textsc{Hansen} im Herbst 1826 an der hiesigen Sternwarte praktisch kennengelernt.
Abgelesen wurde mit Mikroskopen statt mit Nonien. Die folgende Übersicht enthält
die Hauptergebnisse dieser Beobachtungen. Die erste Spalte gibt die Bezeichnung des
Sterns an, die zweite die Lage des Kreises, die dritte die Anzahl der Beobachtungen,
die vierte die von mir auf den Anfang des Jahres 1827 reduzierte Zenitdistanz
(nördliche mit positivem Vorzeichen), die fünfte die Breite, die sich aus den oben
auf S.~9 mitgeteilten Deklinationen ergibt.

% Printed 53; principal PDF 60; Google PDF 64.
\par\clearpage
\noindent\makebox[12mm][l]{}\parbox[t]{\dimexpr\linewidth-24mm\relax}{\centering\scriptsize\bfseries BESTIMMUNG DES BREITENUNTERSCHIEDS ETC.\quad IV. BREITENBESTIMMUNG VON SEEBERG.}\makebox[12mm][r]{\fontsize{13}{16}\selectfont 53}\par\medskip

{\small
\setlength{\arraycolsep}{3.2pt}
\renewcommand{\arraystretch}{1.06}
\[
\begin{array}{@{}r|c|c|r@{\;}r@{\;}r|r||r|c|c|r@{\;}r@{\;}r|r@{}}
1&\text{Ost}&5&-1^\circ&1'&53''{,}14&50^\circ56'4''{,}76&
16&\text{Ost}&5&-1^\circ&42'&19''{,}87&50^\circ56'6''{,}55\\
&\text{West}&4&&&52{,}80&4{,}42&&\text{West}&5&&&17{,}05&3{,}73\\
2&\text{Ost}&5&+4&37&29{,}41&5{,}57&17&\text{Ost}&5&+0&38&48{,}22&5{,}50\\
&\text{West}&6&&&30{,}86&4{,}12&&\text{West}&5&&&48{,}79&4{,}93\\
3&\text{Ost}&5&-0&45&19{,}00&5{,}20&18&\text{Ost}&5&+1&38&59{,}63&6{,}24\\
&\text{West}&7&&&17{,}70&3{,}90&&\text{West}&6&&&61{,}51&4{,}36\\
4&\text{Ost}&6&+3&38&50{,}38&5{,}17&19&\text{West}&2&+3&41&30{,}38&4{,}49\\
&\text{West}&6&&&51{,}05&4{,}50&20&\text{Ost}&1&+3&33&50{,}19&4{,}11\\
5&\text{Ost}&6&+4&29&54{,}35&4{,}87&&\text{West}&1&&&47{,}87&6{,}43\\
&\text{West}&6&&&54{,}30&4{,}92&21&\text{Ost}&1&+1&58&31{,}74&5{,}35\\
6&\text{Ost}&6&-4&20&26{,}87&6{,}01&&\text{West}&1&&&32{,}90&4{,}19\\
&\text{West}&4&&&25{,}92&5{,}06&22&\text{Ost}&1&-4&39&61{,}37&6{,}07\\
7&\text{Ost}&6&-0&39&22{,}12&5{,}56&&\text{West}&1&&&58{,}90&3{,}60\\
&\text{West}&4&&&20{,}70&4{,}14&23&\text{Ost}&1&+5&24&21{,}09&5{,}91\\
8&\text{Ost}&5&+1&40&1{,}41&6{,}06&&\text{West}&1&&&22{,}88&4{,}12\\
&\text{West}&4&&&3{,}24&4{,}23&24&\text{West}&1&-0&33&26{,}54&5{,}95\\
9&\text{Ost}&5&+5&17&29{,}69&6{,}82&25&\text{West}&1&+8&5&40{,}64&6{,}50\\
&\text{West}&4&&&33{,}72&2{,}79&26&\text{West}&1&-0&17&51{,}91&5{,}03\\
10&\text{Ost}&5&+1&43&6{,}51&5{,}54&27&\text{West}&1&-4&35&46{,}53&3{,}57\\
&\text{West}&5&&&8{,}63&3{,}42&28&\text{West}&1&+2&44&9{,}48&3{,}93\\
11&\text{Ost}&5&+2&43&27{,}76&5{,}26&29&\text{West}&1&+1&31&5{,}78&5{,}28\\
&\text{West}&4&&&28{,}89&4{,}13&30&\text{West}&2&+4&40&0{,}22&4{,}51\\
12&\text{Ost}&5&+2&2&48{,}89&6{,}77&31&\text{West}&1&-4&57&58{,}28&4{,}13\\
&\text{West}&3&&&50{,}29&5{,}37&32&\text{West}&1&+2&18&59{,}76&3{,}88\\
13&\text{Ost}&5&-1&29&56{,}73&6{,}19&33&\text{West}&1&-0&39&51{,}61&5{,}28\\
&\text{West}&5&&&56{,}64&6{,}10&34&\text{West}&1&+6&9&34{,}85&2{,}94\\
14&\text{Ost}&5&-0&35&42{,}51&5{,}00&35&\text{West}&1&-3&59&15{,}56&4{,}09\\
&\text{West}&3&&&42{,}84&5{,}33&36&\text{West}&2&-4&6&42{,}83&4{,}84\\
15&\text{Ost}&5&-2&36&14{,}46&6{,}93&37&\text{West}&1&+6&0&36{,}97&5{,}38\\
&\text{West}&4&&&12{,}88&5{,}35&38&\text{West}&1&-1&53&16{,}75&3{,}84
\end{array}
\]
\par
}

Diese sechzig Ergebnisse für die Breite sind freilich unterschiedlich zuverlässig.
Um die ihnen zuzuweisenden Gewichte ohne Willkür angeben zu können, müsste man das
Verhältnis des mittleren eigentlichen Beobachtungsfehlers zum mittleren Teilungsfehler
kennen. Ist dieses Verhältnis wie \(1\) zu \(\surd\theta\), so ist, wenn man die
geringe verbleibende Unsicherheit der Deklinationen außer Acht lässt,
\[
\frac n{1+n\theta}
\]
\noindent das Gewicht einer Bestimmung, die auf \(n\) Beobachtungen am selben
Teilstrich beruht.

Nimmt man statt dieses Gewichts schlechthin \(n\) an, so wird das Mittel
aus den 206 Beobachtungen
\[
=50^\circ56'5''{,}16.
\]

% Printed 54; principal PDF 61; Google PDF 65.
\par\clearpage
\noindent\makebox[12mm][l]{\fontsize{13}{16}\selectfont 54}\parbox[t]{\dimexpr\linewidth-24mm\relax}{\centering\scriptsize\bfseries BESTIMMUNG DES BREITENUNTERSCHIEDS ZWISCHEN GÖTTINGEN UND ALTONA.}\makebox[12mm][r]{}\par\medskip

Die Beobachtungen lassen indessen erkennen, dass die Teilungsfehler beträchtlich
größer sein müssen als am \textsc{Ramsden}schen Zenitsektor, während die eigentlichen
Beobachtungsfehler vielleicht sogar noch etwas kleiner sind. Bei jenem Verfahren
werden die Bestimmungen, die auf einer größeren Anzahl von Beobachtungen beruhen,
gegenüber denen, denen nur eine oder zwei Beobachtungen zugrunde liegen, also
viel zu stark bevorzugt.

Will man jedoch den Einfluss der Teilungsfehler berücksichtigen, muss man auch
beachten, dass jede Bestimmung des Kollimationsfehlers einen konstanten Fehler
enthält, der von den Fehlern der dabei beteiligten Teilstriche abhängt. Dieser
Fehler wirkt sich offensichtlich auf die Polhöhe in entgegengesetzter Richtung
aus, je nachdem, ob sich der Kreis östlich oder westlich befindet. Man muss daher
die Beobachtungen für die verschiedenen Lagen des Kreises voneinander trennen,
aus jeder Reihe unter Anwendung des Gewichts \(\dfrac n{1+n\theta}\) für jede
Bestimmung das Mittel berechnen und zuletzt aus diesen beiden Mitteln das
einfache arithmetische Mittel nehmen.

Da \(\theta\) nicht genau bekannt ist, wurde diese Rechnung unter den drei
Annahmen \(\theta=0\), \(\theta=1\), \(\theta=\infty\) durchgeführt. Daraus
ergab sich für die Polhöhe:
\begin{center}
\renewcommand{\arraystretch}{1.18}%
\begin{tabular}{c|c|c|c}
&\(\theta=0\)&\(\theta=1\)&\(\theta=\infty\)\\
\hline
Kreis Ost & \(50^\circ56'5''{,}75\) & \(50^\circ56'5''{,}69\) & \(50^\circ56'5''{,}71\)\\
Kreis West & \hfill \(4{,}62\) & \hfill \(4{,}65\) & \hfill \(4{,}65\)\\
\hline
Polhöhe & \(50\enspace56\enspace5{,}18\) & \(50\enspace56\enspace5{,}17\) & \(50\enspace56\enspace5{,}18\)
\end{tabular}
\end{center}

Man sieht also, dass sich das erste Ergebnis durch die Berücksichtigung der
strengeren Grundsätze gar nicht merklich ändert und man sich an die Zahl
\(50^\circ56'5''{,}17\) halten kann.

Die Biegung des Fernrohrs ist bei diesen Rechnungen noch nicht berücksichtigt.
Nach Angabe von Herrn \textsc{Hansen} ist sie am Horizont \(=1''{,}00\) und
ist von der beobachteten Zenitdistanz abzuziehen, also in unserer Schreibweise
\(f=-1''{,}00\). Berücksichtigt man diese Biegung, so fällt die aus den nördlich
des Zenits kulminierenden Sternen ermittelte Polhöhe etwas größer aus, die aus
den südlichen kleiner. Da jene etwas über-
\par

% Printed 55; principal PDF 62; Google PDF 66.
\par\clearpage
\noindent\makebox[12mm][l]{}\makebox[\dimexpr\linewidth-24mm\relax][c]{\footnotesize\textbf{IV.\quad BREITENBESTIMMUNG DER STERNWARTE SEEBERG.}}\makebox[12mm][r]{\fontsize{13}{16}\selectfont 55}\par\medskip

\noindent wiegen, erhöht sich das mittlere Ergebnis um \(0''{,}02\). Der zweite
Teil der Biegung, also die Biegung bei vertikaler Stellung, lässt sich, da alle
hier vorkommenden Zenitdistanzen klein sind, als konstante Änderung des
Kollimationsfehlers auffassen. Er fällt daher bei unserem Verfahren ebenso
von selbst weg wie die Teilungsfehler der Teilstriche, die bei der Bestimmung
des Kollimationsfehlers beteiligt sind.

Wir erhalten demnach aus diesen Beobachtungen als endgültigen Wert für die Polhöhe
\[
50^\circ56'5''{,}19.
\]

Die erwähnte trigonometrische Verbindung der beiden Sternwarten ergibt, wenn
man sie mit den oben angegebenen Dimensionen des Erdsphäroids berechnet, den
Breitenunterschied
\[
35'41''{,}86,
\]
\noindent und somit zusammen mit der oben bestimmten Polhöhe von Göttingen die
der Seeberger Sternwarte
\[
=50^\circ56'5''{,}99.
\]

Diese bezieht sich auf den Dreieckspunkt, nämlich den Mittelpunkt der Achse
des Mittagsfernrohrs. Der Mittelpunkt der Achse des Meridiankreises liegt
\(1{,}168\) Toisen oder \(0''{,}07\) im Bogen weiter südlich. Die Polhöhe des
letzteren Punktes ist also, aus Göttingen über die trigonometrische Verbindung abgeleitet,
\[
=50^\circ56'5''{,}92
\]
\noindent also um \(0''{,}73\) größer als nach den astronomischen Beobachtungen.

Für den Längenunterschied ergeben sich aus der trigonometrischen Verbindung
übrigens \(47'9''{,}20\) im Bogen oder \(3^{\mathrm m}8^{\mathrm s}{,}61\) in Zeit,
was sehr gut mit dem aus astronomischen Beobachtungen bekannten Wert übereinstimmt.
Schließlich ergibt sich aus jenen Messungen für das Azimut der Dreiecksseite
Seeberg---südliches Meridianzeichen bei Schwabhausen \(4''{,}6\) westlich.
Auch dies kann als gute Übereinstimmung gelten, wenn man die beträchtliche
Zahl der Zwischenpunkte, die Unterschiede in den Angaben einiger Winkel der
preußischen Messung und die Ungewissheit berücksichtigt, ob der Dreieckspunkt
genau im Meridian lag.

\begin{center}\rule{.24\textwidth}{.4pt}\end{center}

% Printed 56 (unfoliated); principal PDF 63; Google PDF 67.
\par\clearpage

\vspace*{.055\textheight}
\begin{center}{\Large ZUSATZ ZU [ART. 20] S. 48.}\end{center}

\textsc{Walbeck}s Bestimmung der Dimensionen des Erdsphäroids findet sich in einer
kleinen Abhandlung: \textit{De forma et magnitudine telluris, ex dimensis
arcubus meridiani, definiendis}, von der aber nur die ersten beiden Bogen im
Druck erschienen sind (Åbo, 1819). \textsc{Walbeck} hat die peruanische, die beiden
ostindischen, die französische, die englische und die neuere lappländische Gradmessung
in seine Berechnung einbezogen und ist meines Wissens bisher der Einzige, der diese
Aufgabe nach richtigen, von Willkür freien Grundsätzen ausgeführt hat. Er hat jedoch
bei jeder einzelnen Gradmessung nur den ganzen Bogen oder die an den Endpunkten
beobachteten Polhöhen in Betracht gezogen, ohne die bei mehreren Messungen
vorhandenen Zwischenpunkte zu berücksichtigen. In der Rechnung ist er bei der
ersten Potenz der Abplattung stehen geblieben.

Ich habe deshalb Herrn Dr. \textsc{Schmidt}, der sich bereits durch mehrere Arbeiten
einen guten Namen gemacht hatte, vor Kurzem zu einer neuen Berechnung sämtlicher
Gradmessungen veranlasst. Er hat sie während des Drucks der letzten Bogen dieser
Schrift vollendet. Dabei hat er sowohl die höheren Potenzen der Abplattung als
auch die an allen Zwischenpunkten beobachteten Polhöhen berücksichtigt, außerdem
die hannoversche Gradmessung einbezogen und nach dem oben auf S.~50 angedeuteten
Prinzip dasjenige Ellipsoid bestimmt, auf dem die astronomisch beobachteten
Polhöhen möglichst wenig geändert werden müssen, um mit den geodätischen
Messungen vollständig übereinzustimmen, d.\ h. auf dem die Summe der Quadrate
der dazu erforderlichen Änderungen ein Minimum wird. Das Ergebnis dieser
Rechnung ist:
\par

% Printed 57; principal PDF 64; Google PDF 68.
\par\clearpage
\noindent\makebox[12mm][l]{}\makebox[\dimexpr\linewidth-24mm\relax][c]{\footnotesize\textbf{BESTIMMUNG DES BREITENUNTERSCHIEDS ETC.\quad ZUSATZ.}}\makebox[12mm][r]{\fontsize{13}{16}\selectfont 57}\par\medskip

\begin{center}
\begin{tabular}{@{}l@{\quad}l@{}}
\makebox[.62\linewidth][s]{Abplattung\hfill .\hfill .\hfill .\hfill .\hfill .\hfill .\hfill .\hfill .\hfill .\hfill .\hfill .\hfill .\hfill .}
& \(\dfrac1{298{,}39}\)\\
\makebox[.62\linewidth][s]{Dreihundertsechzigster Teil des Erdmeridians\hfill .}
& \(57010{,}35\) Toisen.
\end{tabular}
\end{center}

Die folgende Übersicht zeigt die beobachteten Polhöhen an den 25 Punkten der
sieben Gradmessungen und die kleinsten Änderungen, die für ihre vollständige
Übereinstimmung mit den ermittelten Erddimensionen erforderlich sind:
\[
\renewcommand{\arraystretch}{1.38}
\begin{array}{@{}l@{\quad}|r@{\enspace}r@{\enspace}r|l@{}}
\multicolumn{5}{c}{\textit{Peruanische Messung.}}\\[.3ex]
\text{Tarqui}&-3^\circ&4'&30''{,}83&+2''{,}05\\
\text{Cotchesqui}&+0&2&37{,}83&-2{,}05\\[.8ex]
\multicolumn{5}{c}{\textit{Erste ostindische Messung.}}\\[.3ex]
\text{Trivandeporum}&+11&44&52{,}59&-0{,}48\\
\text{Paudree}&13&19&49{,}02&+0{,}47\\[.8ex]
\multicolumn{5}{c}{\textit{Zweite ostindische Messung.}}\\[.3ex]
\text{Punnae}&8&9&38{,}39&-1{,}43\\
\text{Putchapolliam}&10&59&48{,}93&-1{,}18\\
\text{Dodagoontah}&12&59&59{,}91&+3{,}37\\
\text{Namthabad}&15&6&0{,}64&-0{,}77\\[.8ex]
\multicolumn{5}{c}{\textit{Französische Messung.}}\\[.3ex]
\text{Formentera}&38&39&56{,}11&+3{,}95\\
\text{Montjouy}&41&21&45{,}45&+2{,}81\\
\text{Barcelona}&41&22&47{,}16&+1{,}07\\
\text{Perpignan}&42&41&58{,}01&-3{,}67\\
\text{Carcassonne}&43&12&54{,}31&-0{,}96\\
\text{Evaux}&46&10&42{,}19&-6{,}14\\
\text{Pantheon}&48&50&48{,}94&-0{,}17\\
\text{Dünkirchen}&51&2&8{,}74&+3{,}12\\[.8ex]
\multicolumn{5}{c}{\textit{Englische Messung.}}\\[.3ex]
\text{Dunnose}&50&37&7{,}81&-1{,}73\\
\text{Greenwich}&51&28&39{,}60&+1{,}00\\
\text{Blenheim}&51&50&27{,}50&+3{,}02\\
\text{Arburyhill}&52&13&27{,}79&+1{,}80\\
\text{Clifton}&53&27&31{,}59&-4{,}07
\end{array}
\]
\noindent\textbf{IX.}\hfill 8\hspace{.12\textwidth}\par

% Printed 58; principal PDF 65; Google PDF 69.
\par\clearpage
\noindent\makebox[12mm][l]{\fontsize{13}{16}\selectfont 58}\makebox[\dimexpr\linewidth-24mm\relax][c]{\footnotesize\textbf{BESTIMMUNG DES BREITENUNTERSCHIEDS ETC.\quad ZUSATZ.}}\makebox[12mm][r]{}\par\medskip

\[
\renewcommand{\arraystretch}{1.38}
\begin{array}{@{}l@{\quad}|r@{\enspace}r@{\enspace}r|l@{}}
\multicolumn{5}{c}{\textit{Hannoversche Messung.}}\\[.3ex]
\text{Göttingen}&51^\circ&31'&47''{,}85&-2''{,}65\\
\text{Altona}&53&32&45{,}27&+2{,}66\\[.8ex]
\multicolumn{5}{c}{\textit{Schwedische Messung.}}\\[.3ex]
\text{Mallörn}&65&31&31{,}06&+1{,}40\\
\text{Pahtavara}&67&8&51{,}41&-1{,}40
\end{array}
\]

Die Zahlen in der letzten Spalte sind keineswegs als Fehler der astronomischen
Beobachtungen zu verstehen. Sie sind die algebraische Summe dieser Fehler und
der Unregelmäßigkeiten in der Richtung der Vertikalen. Behandelt man diese
Gesamtabweichungen nach denselben Regeln wie die zufälligen Fehler, so ergibt
sich die mittlere Abweichung \(3''{,}18\) und damit der mittlere zu erwartende Fehler
\begin{center}
\begin{tabular}{@{}l@{\quad}l@{}}
\makebox[.64\linewidth][s]{in dem Nenner der Abplattung\hfill .\hfill .\hfill .\hfill .\hfill .\hfill .}
& \(12{,}5\) Einheiten\\
in dem Wert des dreihundertsechzigsten Teils &\\
\makebox[.64\linewidth][r]{des Erdmeridians} & \(5{,}0\) Toisen.
\end{tabular}
\end{center}

Der sogenannte wahrscheinliche Fehler lässt sich also auf 8 Einheiten beim
Nenner der Abplattung und auf 3 Toisen beim mittleren Breitengrad schätzen.
Damit ist der Grad der Genauigkeit bestimmt, den man der Ermittlung der
Dimensionen des Erdsphäroids durch alle bisherigen Breitengradmessungen
zuschreiben darf. Dies ist als wichtiges Ergebnis dieser verdienstvollen
Arbeit von Herrn Dr. \textsc{Schmidt} anzusehen, die an anderer Stelle ausführlich
bekannt gemacht werden soll.

\begin{center}\rule{.24\textwidth}{.4pt}\end{center}

% Printed 59 (unfoliated); principal PDF 66; Google PDF 70.
\par\clearpage

\vspace*{.075\textheight}
\begin{center}{\LARGE A\,N\,Z\,E\,I\,G\,E.}\end{center}
\vspace{.045\textheight}
\begin{center}
\rule{.52\textwidth}{.4pt}\\[-.2ex]
{\small Göttingische gelehrte Anzeigen.\enspace 16. Juni 1828.}\\[-.2ex]
\rule{.52\textwidth}{.4pt}
\end{center}

Bei Vandenhoeck und Ruprecht: \textit{Bestimmung des Breitenunterschiedes
zwischen den Sternwarten von Göttingen und Altona durch Beobachtungen am}
\textsc{Ramsden}\textit{schen Zenithsector, von} \textsc{Carl Friedrich Gauss}.
1828. 84 Seiten in 4.

Die Messungen, die der Verfasser dieser Schrift in den vergangenen Jahren im
Königreich Hannover ausgeführt hat, hatten zunächst den Zweck, die von Herrn
Professor \textsc{Schumacher} in den dänischen Staaten unternommene Gradmessung um
zwei Grad weiter nach Süden auszudehnen. Eine Dreieckskette von der südlichen
Grenze des Königreichs Hannover bis Hamburg wurde in den Jahren 1821---1823
vollendet und mit dem dänischen Dreieckssystem verbunden. Um daraus ein
eigenständiges Ergebnis als Gradmessung gewinnen zu können, war noch die
astronomische Bestimmung der Krümmung des ganzen Meridianbogens erforderlich.
Sie musste jedoch bis zum Jahr 1827 zurückgestellt werden, weil das
Dreieckssystem in den folgenden Jahren auf höheren Befehl erweitert wurde,
um es an die \textsc{Krayenhoff}schen Dreiecke anzuschließen. Da die Veröffentlichung
der trigonometrischen Arbeiten aus mehreren hier nicht anzuführenden Gründen
einem größeren Werk vorbehalten bleibt, hat Hofrat \textsc{Gauss} den astronomischen
Teil der Gradmessung nun gleich als eigenständiges Werk herausgegeben, damit
das Endergebnis sofort den übrigen Gradmessungen beigefügt und verwendet
werden kann.

Günstiger gelegene Endpunkte als diejenigen, die sich hier von selbst anboten,
hätte man sich gar nicht wünschen können. Es sind die
\par\smallskip
\noindent\hfill \(8^{*}\)\hspace{.12\textwidth}\par

% Printed 60; principal PDF 67; Google PDF 71.
\par\clearpage
\noindent\makebox[12mm][l]{\fontsize{13}{16}\selectfont 60}\makebox[\dimexpr\linewidth-24mm\relax][c]{\small\textbf{ANZEIGE.}}\makebox[12mm][r]{}\par\medskip

\noindent Sternwarten von Göttingen und Altona, beide mit ausgezeichneten
Instrumenten ausgerüstet und beide durch ein einzigartiges Spiel des Zufalls
so genau auf ein und demselben Meridian gelegen, dass man bestimmte Orte in
den Sternwarten angeben muss, um einen Unterschied feststellen zu können.
Bezogen auf die Mittelpunkte der Achsen der \textsc{Reichenbach}schen Meridiankreise
liegt nämlich die Altonaer Sternwarte nur \(7\tfrac12\) Toisen weiter westlich.

Die Ausdehnung des Himmelsbogens, der dem Stück des Erdmeridians zwischen diesen
beiden Sternwarten entspricht, ergibt sich zwar schon aus ihren absoluten
Polhöhen, zu deren Bestimmung die fortgesetzten Beobachtungen immer neue Beiträge
liefern. Dennoch war es von großer Bedeutung, den Breitenunterschied noch auf
eine andere Weise mit ein und demselben erstklassigen Instrument zu bestimmen.
Hofrat \textsc{Gauss} konnte dafür den ausgezeichneten \textsc{Ramsden}schen
Zenitsektor benutzen, der für ähnliche Arbeiten bei der englischen Gradmessung
verwendet und bekanntlich von \textsc{Mudge} ausführlich beschrieben wurde.

Das Werk gliedert sich in vier Abschnitte. Im ersten werden die beobachteten
Sterne aufgeführt. Charakteristisch für die neuere Beobachtungskunst ist, dass
bei wichtigen Anlässen die Zahl der Beobachtungen stark erhöht wird. So lässt
sich der Einfluss der Unvollkommenheit der Sinne und der Instrumente (denn
absolut vollkommen \textls[100]{kann} keines sein) sowie der unvermeidlichen
äußeren Störquellen auf einen sehr kleinen Teil seiner sonstigen Größe
herabsetzen, auch wenn es unmöglich ist, ihn ganz zu beseitigen. Dieses Ziel
lässt sich nur erreichen, indem man die einzelnen Fehlerquellen bei einer
großen Zahl von Beobachtungen auf vielfältige Weise zur Wirkung kommen lässt.
Insbesondere können daher die unregelmäßigen Teilungsfehler eines nicht
wiederholenden Instruments nur durch die Beobachtung einer beträchtlichen
Anzahl von Sternen mit unterschiedlichen Deklinationen möglichst gut
ausgeglichen werden. Der Verfasser wählte deshalb 43 Sterne in geeigneten
Lagen zur Beobachtung aus, von denen einige noch in keinem Sternverzeichnis
aufgeführt sind.

Der zweite Abschnitt behandelt die Beobachtungen selbst, die in Göttingen
von Anfang April bis Mitte Mai und danach in Altona im Juni angestellt wurden
und deren Anzahl sich auf 900
\par

% Printed 61; principal PDF 68; Google PDF 72.
\par\clearpage
\noindent\makebox[12mm][l]{}\parbox[t]{\dimexpr\linewidth-24mm\relax}{\centering\scriptsize\bfseries BESTIMMUNG DES BREITENUNTERSCHIEDS ZWISCHEN GÖTTINGEN UND ALTONA.}\makebox[12mm][r]{\fontsize{13}{16}\selectfont 61}\par\medskip

\noindent beläuft. Sie sind alle einzeln aufgeführt, gleich nach Sternen
geordnet und mit ihren Reduktionen auf die mittlere Stellung für den Anfang des
Jahres 1827. Die ursprüngliche Form des Tagebuchs erschien dem Verfasser
unnötig und hätte den Umfang des Werks verdoppelt.

Der dritte Abschnitt entwickelt die Ergebnisse aus den Beobachtungen. Das
Hauptergebnis ist der Breitenunterschied zwischen den beiden Sternwarten, der
für die Standorte der Meridiankreise nach mehreren von verschiedenen
Gesichtspunkten ausgehenden Kombinationen auf \(2^\circ0'57''{,}42\) festgesetzt
wird. Dabei bleibt nur eine wahrscheinliche Unsicherheit von \(0''{,}07\).
Bemerkenswert ist die aus den Beobachtungen ganz deutlich hervorgehende
Änderung des Kollimationsfehlers. Dieser betrug bei den Beobachtungen in
Göttingen \(3''{,}76\), bei denen in Altona \(1''{,}39\) und muss sich während
des mit größter Vorsicht durchgeführten Transports geändert haben. Diese
Tatsache ist besonders wichtig für die Einschätzung des Vertrauens, das die
ältere lappländische Gradmessung verdient. Ihr ist mit Recht vorgeworfen
worden, dass der Zenitsektor an den einzelnen Beobachtungsorten nicht
umgewendet wurde. Dass der Kollimationsfehler bei ihr gleich blieb, wird durch
das, was dazu neuerdings in einer sonst sehr schätzenswerten Abhandlung
vorgebracht wurde, keineswegs zufriedenstellend begründet. Man findet hier
ferner die Beobachtungen des Verfassers aus den Jahren 1820 und 1824, die zur
Bestimmung der absoluten Polhöhe der Göttinger Sternwarte dienen, den Vergleich
des Ergebnisses für den Breitengrad aus der hannoverschen Gradmessung mit
den Dimensionen des Erdsphäroids, die \textsc{Walbeck} 1819 aus sämtlichen
zuverlässigen früheren Gradmessungen abgeleitet hat, die Ansichten des
Verfassers zu den Unregelmäßigkeiten der Erdfigur überhaupt und schließlich
den Vergleich der Deklinationen der Zenitsterne aus den vorliegenden
Beobachtungen mit den Bestimmungen von \textsc{Bradley} und \textsc{Piazzi}.
Bemerkenswert ist bei einem der kleineren Sterne (\textsc{Piazzi} 16. 291)
eine Differenz von \(10''{,}3\) gegenüber \textsc{Piazzi}s Katalog. Nach den vom
Verfasser vorgebrachten Gründen ist darin eine Eigenbewegung unverkennbar.

Der vierte Abschnitt enthält die Bestimmung der Breite der Sternwarte
Seeberg aus den Beobachtungen, die der Leiter dieser Sternwarte,
\par
\par
\endgroup
